% !TEX program = xelatex
\documentclass[aspectratio=169,11pt,handout]{beamer}
\usetheme{Madrid}
\usepackage{fontspec}
\setsansfont{DejaVu Sans}
\setmainfont{DejaVu Serif}
\usepackage{polyglossia}
\setdefaultlanguage{turkish}
\usepackage{amsmath,amssymb}
\usepackage{booktabs}
\usepackage{multicol}
\usepackage{xcolor}

% ---------- Renkler ----------
\definecolor{anamavi}{RGB}{0,45,98}
\definecolor{cozumyesil}{RGB}{0,110,50}
\definecolor{uyarikirmizi}{RGB}{180,20,20}
\setbeamercolor{structure}{fg=anamavi}
\setbeamercolor{frametitle}{bg=anamavi,fg=white}
\setbeamercolor{title}{bg=anamavi,fg=white}
\setbeamercolor{palette primary}{bg=anamavi,fg=white}
\setbeamercolor{palette secondary}{bg=anamavi!85!black,fg=white}
\setbeamercolor{palette tertiary}{bg=anamavi!70!black,fg=white}
\setbeamercolor{block title}{bg=anamavi,fg=white}
\setbeamercolor{block body}{bg=anamavi!7}
\setbeamercolor{block title example}{bg=cozumyesil,fg=white}
\setbeamercolor{block body example}{bg=cozumyesil!8}
\setbeamercolor{block title alerted}{bg=uyarikirmizi,fg=white}
\setbeamercolor{block body alerted}{bg=uyarikirmizi!7}
\setbeamertemplate{navigation symbols}{}
\setbeamertemplate{blocks}[rounded][shadow=false]

% ---------- Soru / Çözüm kalıbı ----------
% \soru{başlık}{soru metni}{çözüm metni}: önce soru, bir tıklamayla çözüm
\newcommand{\soru}[3]{%
\begin{frame}{#1}
  \begin{block}{Soru}\large #2\end{block}
  \pause
  \begin{exampleblock}{Çözüm}\large #3\end{exampleblock}
\end{frame}}
\newcommand{\ab}[1]{\overline{#1}}

\title[MMAT105 -- Hafta 13]{MMAT105 Genel Matematik}
\subtitle{Hafta 13: Sayı ve Kesir Problemleri, Yaş Problemleri}
\author{}
\date{2026--2027 Güz Dönemi}

\begin{document}

% =====================================================
\begin{frame}
\titlepage
\end{frame}

% (iletişim çerçevesi kaldırıldı)

\begin{frame}{Geçen Hafta Hatırlatma}
\begin{columns}[t]
\column{0.58\textwidth}
\textbf{Hesaplayınız:}
\begin{enumerate}
\item \(\dfrac{a}{b}=\dfrac{3}{5}\) ve \(a+b=48\) ise \(a=?\)
\item \(4,\ 7,\ 10,\ 15\) sayılarının aritmetik ortalaması?
\item \(x\) ile \(y\) doğru orantılıdır. \(x=4\) iken \(y=10\) ise \(x=6\) iken \(y=?\)
\end{enumerate}
\column{0.36\textwidth}
\pause
\textbf{Cevaplar:}
\begin{enumerate}
\item \(a=3k,\ b=5k\), \(8k=48\Rightarrow a=18\)
\item \(\dfrac{36}{4}=9\)
\item \(\dfrac{y}{x}=\dfrac{10}{4}\Rightarrow y=15\)
\end{enumerate}
\end{columns}
\end{frame}

\begin{frame}{Bugün}
\begin{columns}[t]
\column{0.32\textwidth}
\begin{block}{Ders 1}
\begin{enumerate}
\item Sözel ifade \(\to\) denklem
\item Toplam, fark, oran
\item Ardışık sayılar
\item Basamak problemleri
\item Sıra ve madeni para
\end{enumerate}
\end{block}
\column{0.32\textwidth}
\begin{block}{Ders 2}
\begin{enumerate}\setcounter{enumi}{5}
\item Bütünün kesri
\item Kalanın kesri
\item EKOK kadar birim alma
\item Depo doluluk oranı
\item Kesirli denklemler
\end{enumerate}
\end{block}
\column{0.32\textwidth}
\begin{block}{Ders 3}
\begin{enumerate}\setcounter{enumi}{10}
\item Yaş farkı değişmez
\item Toplam \(n\cdot t\) artar
\item Tablo yöntemi
\item ``Ben senin yaşındayken''
\item Yaş oranı, doğum yılı
\end{enumerate}
\end{block}
\end{columns}
\end{frame}

% =====================================================
\section{Ders 1: Sayı Problemleri}
\begin{frame}
\centering\Large\color{anamavi}\textbf{Ders 1}\\[6pt]\normalsize Sözel İfadeyi Denkleme Çevirme ve Sayı Problemleri
\end{frame}

\subsection{Denklem kurma}
\begin{frame}{Sözel İfadeden Denkleme (sayı \(=x\))}
\begin{center}
\renewcommand{\arraystretch}{1.35}
\begin{tabular}{ll}
\toprule
Sözel ifade & Cebirsel ifade\\\midrule
Bir sayının \(3\) katının \(5\) eksiği & \(3x-5\)\\
Bir sayının \(5\) eksiğinin \(3\) katı & \(3(x-5)\)\\
Bir sayının \(2\) katından \(7\) fazlası & \(2x+7\)\\
Bir sayının yarısı ile üçte birinin toplamı & \(\tfrac{x}{2}+\tfrac{x}{3}\)\\
Bir sayının \(4\) fazlasının yarısı & \(\tfrac{x+4}{2}\)\\
\(a\), \(b\)'nin \(3\) katından \(2\) eksiktir & \(a=3b-2\)\\\bottomrule
\end{tabular}
\end{center}
\begin{alertblock}{Sık yapılan hatalar}
``\(5\) eksiğinin \(3\) katı'' \(=3(x-5)\), \ \(3x-5\) \textbf{değildir}: önce yapılan işlem paranteze girer.
\end{alertblock}
\end{frame}

\begin{frame}{Problem Çözme Adımları}
\begin{block}{Dört adım}
\begin{enumerate}
\item Bilinmeyeni seç ve adlandır (\(x\) \(=\) ``küçük sayı'', ``sıra sayısı''\dots).
\item Diğer büyüklükleri \(x\) cinsinden yaz.
\item Verilen eşitliği denkleme çevir ve çöz (5.--6. hafta).
\item \textbf{Sağlama:} bulunan değeri sorunun metnine yerleştir; cevap anlamlı mı? (tam sayı, pozitif, rakam \(0\)--\(9\) arası\dots)
\end{enumerate}
\end{block}
\begin{alertblock}{Sık yapılan hatalar}
Bulunan \(x\) her zaman istenen cevap değildir: soru ``büyük sayı'', ``en büyüğü'', ``sınıf mevcudu'' diye sorabilir.
\end{alertblock}
\end{frame}

\soru{Denklem Kurma -- Örnek 1}
{Hangi sayının \(3\) katının \(5\) eksiği, aynı sayının \(2\) katından \(7\) fazlasına eşittir?}
{\[3x-5=2x+7\ \Rightarrow\ x=12\]
Sağlama: \(3\cdot 12-5=31\), \quad \(2\cdot 12+7=31\) \checkmark}

\soru{Denklem Kurma -- Örnek 2}
{Bir sayının yarısı ile üçte birinin toplamı \(35\)'tir. Bu sayı kaçtır?}
{Paydaların EKOK'u \(6\): denklemi \(6\) ile çarpalım.
\[\frac{x}{2}+\frac{x}{3}=35\ \Rightarrow\ 3x+2x=210\ \Rightarrow\ x=42\]
Sağlama: \(21+14=35\) \checkmark}

\soru{Denklem Kurma -- Örnek 3}
{Bir sayının \(4\) fazlasının \(3\) katı, aynı sayının \(5\) katından \(6\) eksiktir. Bu sayı kaçtır?}
{\[3(x+4)=5x-6\ \Rightarrow\ 3x+12=5x-6\ \Rightarrow\ 2x=18\ \Rightarrow\ x=9\]
Sağlama: \(3\cdot 13=39\), \quad \(5\cdot 9-6=39\) \checkmark}

\subsection{Toplam, fark, oran}
\soru{İki Sayı -- Örnek 4}
{İki sayının toplamı \(64\), farkı \(18\)'dir. Bu sayıları bulunuz.}
{\(x+y=64\) ve \(x-y=18\). Taraf tarafa toplayalım:
\[2x=82\ \Rightarrow\ x=41,\qquad y=64-41=23\]
Kısa yol: büyük \(=\dfrac{\text{toplam}+\text{fark}}{2}\), \ küçük \(=\dfrac{\text{toplam}-\text{fark}}{2}\).}

\soru{İki Sayı -- Örnek 5}
{İki sayının oranı \(\dfrac{3}{5}\), toplamı \(96\)'dır. Bu sayılar kaçtır?}
{Sayılar \(3k\) ve \(5k\) (12. hafta):
\[3k+5k=96\ \Rightarrow\ k=12\ \Rightarrow\ 36\ \text{ve}\ 60\]
Sağlama: \(36+60=96\), \ \(\dfrac{36}{60}=\dfrac{3}{5}\) \checkmark}

\soru{İki Sayı -- Örnek 6}
{Toplamları \(70\) olan iki doğal sayıdan büyüğü küçüğüne bölündüğünde bölüm \(4\), kalan \(5\) oluyor. Büyük sayı kaçtır?}
{Küçük \(=k\) \(\Rightarrow\) büyük \(=4k+5\) (2. hafta: \(A=B\cdot C+K\)).
\[k+4k+5=70\ \Rightarrow\ 5k=65\ \Rightarrow\ k=13,\quad \text{büyük}=57\]
Sağlama: \(57=13\cdot 4+5\) ve \(5<13\) \checkmark}

\subsection{Ardışık sayılar}
\begin{frame}{Ardışık Sayılar}
\begin{block}{Kural}
\begin{itemize}
\item Ardışık tam sayılar: \(x,\ x+1,\ x+2,\ \dots\)
\item Ardışık çift / tek sayılar: \(x,\ x+2,\ x+4,\ \dots\) \ (\(x\) çift / tek)
\item Tek sayıda ardışık terim: toplam \(=\) terim sayısı \(\times\) ortanca terim
\end{itemize}
\end{block}
\begin{alertblock}{Sık yapılan hatalar}
Ardışık tek sayılar \(x,\ x+1,\ x+2\) \textbf{değildir}; ardışık tek (ve çift) sayılar arasındaki fark \(2\)'dir.
\end{alertblock}
\end{frame}

\soru{Ardışık Sayılar -- Örnek 7}
{Ardışık üç tam sayının toplamı \(87\)'dir. En büyüğü kaçtır?}
{\[x+(x+1)+(x+2)=87\ \Rightarrow\ 3x+3=87\ \Rightarrow\ x=28\]
Sayılar \(28,\ 29,\ 30\) \(\Rightarrow\) en büyüğü \(30\).\\[3pt]
Kısa yol: ortanca \(=87\div 3=29\).}

\soru{Ardışık Sayılar -- Örnek 8}
{Ardışık dört çift sayının toplamı \(100\)'dür. En küçüğü kaçtır?}
{\[x+(x+2)+(x+4)+(x+6)=100\ \Rightarrow\ 4x+12=100\ \Rightarrow\ x=22\]
Sayılar \(22,\ 24,\ 26,\ 28\) \quad (toplam \(100\) \checkmark). \ En küçüğü \(22\).}

\subsection{Basamak problemleri}
\begin{frame}{Basamak Problemleri}
\begin{block}{Kural (\(a,b\) rakam, \(a\neq 0\))}
\[\ab{ab}=10a+b,\qquad \ab{ab}+\ab{ba}=11(a+b),\qquad \ab{ab}-\ab{ba}=9(a-b)\]
Bir sayının sağına \(c\) rakamı yazılırsa yeni sayı \(10\cdot(\text{sayı})+c\) olur.
\end{block}
\begin{alertblock}{Sık yapılan hatalar}
\begin{itemize}
\item \(\ab{ab}\neq a\cdot b\): \ \(\ab{47}=40+7\), \ \(4\cdot 7\) değil.
\item Rakamlar \(0\)--\(9\) arasındadır; \(\ab{ba}\) iki basamaklı olacaksa \(b\neq 0\).
\end{itemize}
\end{alertblock}
\end{frame}

\soru{Basamak -- Örnek 9}
{İki basamaklı bir sayının rakamları toplamı \(11\)'dir. Rakamlarının yerleri değiştirilince sayı \(27\) artıyor. Bu sayı kaçtır?}
{\(\ab{ba}-\ab{ab}=9(b-a)=27\ \Rightarrow\ b-a=3\).
\[a+b=11,\ \ b-a=3\ \Rightarrow\ b=7,\ a=4\ \Rightarrow\ \ab{ab}=47\]
Sağlama: \(74-47=27\) \checkmark}

\soru{Basamak -- Örnek 10}
{İki basamaklı kaç sayı, rakamları toplamının \(4\) katına eşittir?}
{\[10a+b=4(a+b)\ \Rightarrow\ 6a=3b\ \Rightarrow\ b=2a\]
\(b\le 9\) \(\Rightarrow\) \(a=1,2,3,4\): \quad \(12,\ 24,\ 36,\ 48\) \(\Rightarrow\) \(4\) tane.\\[3pt]
Sağlama: \(36=4\cdot(3+6)\) \checkmark}

\subsection{Sıra ve madeni para}
\soru{Sıra Problemi -- Örnek 11}
{Bir sınıftaki öğrenciler sıralara \(2\)'şerli oturursa \(5\) öğrenci ayakta kalıyor; \(3\)'erli oturursa \(2\) sıra boş kalıyor (diğer sıralar tam dolu). Sınıf mevcudu kaçtır?}
{Sıra sayısı \(x\): \quad mevcut \(=2x+5\) \quad ve \quad mevcut \(=3(x-2)\).
\[2x+5=3x-6\ \Rightarrow\ x=11,\qquad \text{mevcut}=2\cdot 11+5=27\]
Sağlama: \(3\cdot 9=27\) \checkmark}

\soru{İki Bilinmeyen -- Örnek 12}
{Bir kumbarada yalnızca \(1\) TL'lik ve \(5\) TL'lik madeni paralar vardır. Toplam \(20\) madeni para ve \(68\) TL olduğuna göre kaç tane \(5\) TL'lik vardır?}
{\(1\) TL'lik \(x\), \(5\) TL'lik \(y\) tane: \quad \(x+y=20\), \ \ \(x+5y=68\).\\[3pt]
Çıkaralım: \(4y=48\ \Rightarrow\ y=12\), \ \(x=8\).\\[3pt]
Sağlama: \(8+60=68\) TL \checkmark}

\begin{frame}{Ders 1 -- Sınıfta Siz Çözün}
\begin{columns}[t]
\column{0.62\textwidth}
\begin{enumerate}
\item Bir sayının \(5\) katının \(8\) eksiği, \(3\) katının \(12\) fazlasına eşit. Sayı? \(\star\)
\item Ardışık üç tek sayının toplamı \(75\). En büyüğü?
\item İki sayının toplamı \(50\), büyüğü küçüğünün \(4\) katı. Küçük sayı?
\item \(\ab{ab}+\ab{ba}=121\) ve \(a-b=5\) ise \(\ab{ab}=?\) \(\star\)
\item Öğrenciler \(3\)'erli oturursa \(4\) kişi ayakta kalıyor, \(4\)'erli oturursa \(3\) sıra boş kalıyor. Mevcut? \(\star\)
\end{enumerate}
\column{0.34\textwidth}
\pause
\textbf{Cevaplar:}
\begin{enumerate}
\item \(10\)
\item \(27\) \ (\(23,25,27\))
\item \(10\)
\item \(83\)
\item \(52\)
\end{enumerate}
\end{columns}
\vspace{4pt}{\footnotesize \(\star\) işaretli soruların adım adım çözümü sonraki slaytlardadır.}
\end{frame}

% ---- Ders 1 alıştırma çözümleri ----
\soru{Ders 1 -- Alıştırma Çözümü (Soru 1)}
{Bir sayının \(5\) katının \(8\) eksiği, aynı sayının \(3\) katının \(12\) fazlasına eşittir. Sayı kaçtır?}
{\(5\) katının \(8\) eksiği: \(5x-8\); \quad \(3\) katının \(12\) fazlası: \(3x+12\).
\[5x-8=3x+12\ \Rightarrow\ 2x=20\ \Rightarrow\ x=10\]
Sağlama: \(50-8=42=30+12\) \checkmark}

\soru{Ders 1 -- Alıştırma Çözümü (Soru 4)}
{\(\ab{ab}+\ab{ba}=121\) ve \(a-b=5\) ise \(\ab{ab}\) kaçtır?}
{\(\ab{ab}+\ab{ba}=(10a+b)+(10b+a)=11(a+b)=121\ \Rightarrow\ a+b=11\).
\[a+b=11,\ \ a-b=5\ \Rightarrow\ 2a=16\ \Rightarrow\ a=8,\ b=3\]
\(\ab{ab}=83\). \quad Sağlama: \(83+38=121\) \checkmark}

\soru{Ders 1 -- Alıştırma Çözümü (Soru 5)}
{Öğrenciler sıralara \(3\)'erli oturursa \(4\) kişi ayakta kalıyor, \(4\)'erli oturursa \(3\) sıra boş kalıyor (diğerleri tam dolu). Mevcut kaçtır?}
{Sıra sayısı \(x\). \quad \(3\)'erli: \(3x+4\); \quad \(4\)'erli: \(x-3\) sıra dolu \(\Rightarrow\) \(4(x-3)\).
\[3x+4=4x-12\ \Rightarrow\ x=16\]
Mevcut \(=3\cdot 16+4=52\). \quad Sağlama: \(4\cdot 13=52\) \checkmark}


% =====================================================
\section{Ders 2: Kesir Problemleri}
\begin{frame}
\centering\Large\color{anamavi}\textbf{Ders 2}\\[6pt]\normalsize Kesir Problemleri: Bütünün Kesri, Kalanın Kesri, Depo Doluluk Oranı
\end{frame}

\subsection{Bütünün ve kalanın kesri}
\begin{frame}{Kesir Problemleri}
\begin{block}{Kural}
\begin{itemize}
\item Bir \(A\) büyüklüğünün \(\dfrac{m}{n}\)'si: \ \(\dfrac{m}{n}\cdot A\). \quad \(\dfrac{m}{n}\)'si \(B\) ise \ \(A=B\cdot\dfrac{n}{m}\).
\item \textbf{Kalanın} kesri: önce kalanı bul, kesri kalana uygula.
\item \textbf{Pratik yol:} bütüne \(x\) yerine paydaların EKOK'u kadar birim (\(12k\), \(15k\) \dots) ver; tüm parçalar tam sayı çıkar.
\end{itemize}
\end{block}
\begin{alertblock}{Sık yapılan hatalar}
``\(\frac{1}{5}\)'ini, sonra \(\frac{1}{3}\)'ünü'' (bütünün) ile ``\(\frac{1}{5}\)'ini, sonra \textbf{kalanın} \(\frac{1}{3}\)'ünü'' farklı sorulardır.
\end{alertblock}
\end{frame}

\soru{Bütünün Kesri -- Örnek 13}
{Bir sayının \(\dfrac{3}{4}\)'ü \(36\)'dır. Bu sayının \(\dfrac{5}{6}\)'sı kaçtır?}
{\[\frac{3}{4}x=36\ \Rightarrow\ x=36\cdot\frac{4}{3}=48\]
\[\frac{5}{6}\cdot 48=40\]}

\soru{Kalanın Kesri -- Örnek 14}
{Ali parasının \(\dfrac{2}{5}\)'ini harcadı, kalanın \(\dfrac{1}{3}\)'ünü kardeşine verdi. Geriye \(120\) TL kaldı. Ali'nin başlangıçta kaç TL'si vardı?}
{Paydalar \(5\) ve \(3\) \(\Rightarrow\) para \(=15k\).\\[2pt]
Harcadı: \(6k\) \(\to\) kalan \(9k\); \quad verdi: \(\frac{1}{3}\cdot 9k=3k\) \(\to\) kalan \(6k\).
\[6k=120\ \Rightarrow\ k=20\ \Rightarrow\ 15k=300\ \text{TL}\]
Sağlama: \(300\to 180\to 120\) \checkmark}

\soru{Bütünün Kesri -- Örnek 15}
{Bir memur maaşının \(\dfrac{1}{5}\)'ini kiraya, \(\dfrac{1}{3}\)'ünü mutfak masraflarına harcıyor. Geriye \(1400\) TL kalıyor. Maaşı kaç TL'dir?}
{Her iki kesir de \textbf{maaşın} kesri. Maaş \(=15k\):
\[\text{kira } 3k,\ \ \text{mutfak } 5k\ \Rightarrow\ \text{kalan } 15k-8k=7k\]
\(7k=1400\ \Rightarrow\ k=200\ \Rightarrow\ \) maaş \(3000\) TL.\\[3pt]
Sağlama: \(3000-600-1000=1400\) \checkmark}

\soru{Kalanın Kesri -- Örnek 16}
{Bir öğrenci kitabın birinci gün \(\dfrac{1}{4}\)'ünü, ikinci gün kalan kısmın \(\dfrac{2}{5}\)'ini okuyor. Geriye \(90\) sayfa kalıyor. Kitap kaç sayfadır?}
{Kitap \(=20k\) \ (\(4\) ve \(5\)'in EKOK'u).\\[2pt]
1. gün: \(5k\) \(\to\) kalan \(15k\); \quad 2. gün: \(\frac{2}{5}\cdot 15k=6k\) \(\to\) kalan \(9k\).
\[9k=90\ \Rightarrow\ k=10\ \Rightarrow\ 200\ \text{sayfa}\]
Sağlama: \(200\to 150\to 90\) \checkmark}

\subsection{Depo doluluk oranı}
\begin{frame}{Depo / Havuz Doluluk Oranı}
\begin{block}{Kural}
Deponun tamamı \(V\) litre olsun. Doluluk \(\dfrac{a}{b}\)'den \(\dfrac{c}{d}\)'ye çıkıyorsa eklenen su:
\[\left(\frac{c}{d}-\frac{a}{b}\right)V\]
Kesir farkı ile eklenen (veya kullanılan) litre eşitlenir.
\end{block}
\begin{alertblock}{Sık yapılan hatalar}
Kesirleri değil, kesirlerin \textbf{farkını} eklenen miktara eşitleyin. Cevap ``deponun tamamı'' mı, ``içindeki su'' mu? Soruyu tekrar okuyun.
\end{alertblock}
\end{frame}

\soru{Depo -- Örnek 17}
{Bir deponun \(\dfrac{3}{4}\)'ü doludur. Depoya \(20\) litre su eklenince deponun \(\dfrac{5}{6}\)'sı doluyor. Deponun tamamı kaç litre su alır?}
{\[\frac{5}{6}-\frac{3}{4}=\frac{10-9}{12}=\frac{1}{12}\ \Rightarrow\ \frac{1}{12}V=20\ \Rightarrow\ V=240\ \text{L}\]
Sağlama: \(180+20=200=\frac{5}{6}\cdot 240\) \checkmark}

\soru{Depo -- Örnek 18}
{Bir su deposunun \(\dfrac{2}{3}\)'ü doludur. Depodan \(30\) litre su kullanılınca deponun yarısı dolu kalıyor. Deponun tamamı kaç litredir?}
{\[\frac{2}{3}-\frac{1}{2}=\frac{1}{6}\ \Rightarrow\ \frac{1}{6}V=30\ \Rightarrow\ V=180\ \text{L}\]
Sağlama: \(120-30=90=\frac{1}{2}\cdot 180\) \checkmark}

\subsection{Kesirli denklemler}
\soru{Kesir -- Örnek 19}
{Bir sınıfın \(\dfrac{3}{7}\)'si kızdır. Sınıfa \(4\) kız öğrenci katılınca kızlar sınıfın yarısı oluyor. Başlangıçta sınıf mevcudu kaçtır?}
{Mevcut \(7k\): \ kız \(3k\), erkek \(4k\). Erkek sayısı değişmez.\\[2pt]
Yeni durumda kız \(=\) erkek: \quad \(3k+4=4k\ \Rightarrow\ k=4\) \(\Rightarrow\) mevcut \(28\).\\[3pt]
Sağlama: \(12+4=16\), \ \(16\) kız \(+\ 16\) erkek \checkmark}

\soru{Kesir -- Örnek 20}
{Bir otobüsteki yolcuların \(\dfrac{1}{3}\)'ü ilk durakta iniyor, \(6\) yolcu biniyor. Yolcu sayısı başlangıçtakinin \(\dfrac{3}{4}\)'ü oluyor. Başlangıçta kaç yolcu vardı?}
{Yolcu \(=12k\): \quad \(12k-4k+6=\dfrac{3}{4}\cdot 12k=9k\)
\[8k+6=9k\ \Rightarrow\ k=6\ \Rightarrow\ 72\ \text{yolcu}\]
Sağlama: \(72-24+6=54=\frac{3}{4}\cdot 72\) \checkmark}

\soru{Kesirli Denklem -- Örnek 21}
{Bir sayının \(\dfrac{2}{3}\)'ünün \(4\) fazlası, aynı sayının \(\dfrac{3}{4}\)'ünün \(1\) eksiğine eşittir. Bu sayı kaçtır?}
{Paydaların EKOK'u \(12\): her terimi \(12\) ile çarpalım.
\[\frac{2x}{3}+4=\frac{3x}{4}-1\ \Rightarrow\ 8x+48=9x-12\ \Rightarrow\ x=60\]
Sağlama: \(40+4=44=45-1\) \checkmark}

\soru{Kesirli Denklem -- Örnek 22}
{Bir kesrin paydası payından \(5\) fazladır. Kesrin payına ve paydasına \(3\) eklenince kesir \(\dfrac{2}{3}\)'e eşit oluyor. Bu kesir nedir?}
{Kesir \(\dfrac{p}{p+5}\):
\[\frac{p+3}{p+8}=\frac{2}{3}\ \Rightarrow\ 3p+9=2p+16\ \Rightarrow\ p=7\]
Kesir \(\dfrac{7}{12}\). \quad Sağlama: \(\dfrac{10}{15}=\dfrac{2}{3}\) \checkmark}

\soru{Geriye Doğru -- Örnek 23}
{Bir top, bırakıldığı yüksekliğin \(\dfrac{2}{5}\)'i kadar zıplamaktadır. Top yere ikinci vuruşundan sonra \(8\) cm yükseldiğine göre ilk bırakıldığı yükseklik kaç cm'dir?}
{Sondan başa (ters işlem: \(\frac{2}{5}\)'i verilmişse \(\cdot\frac{5}{2}\)):
\[8\cdot\frac{5}{2}=20\ \text{cm (1. zıplama)},\qquad 20\cdot\frac{5}{2}=50\ \text{cm}\]
Sağlama: \(50\to 20\to 8\) \checkmark}

\begin{frame}{Ders 2 -- Sınıfta Siz Çözün}
\begin{columns}[t]
\column{0.62\textwidth}
\begin{enumerate}
\item Bir sayının \(\frac{2}{7}\)'si \(18\) ise \(\frac{3}{7}\)'si kaçtır?
\item Parasının \(\frac{1}{3}\)'ünü, sonra kalanın \(\frac{1}{4}\)'ünü harcayan birine \(60\) TL kalıyor. Parası? \(\star\)
\item Deponun \(\frac{2}{5}\)'i dolu; \(60\) L eklenince \(\frac{2}{3}\)'ü doluyor. Depo kaç L? \(\star\)
\item Maaşın \(\frac{1}{4}\)'ü kira, \(\frac{1}{6}\)'sı fatura; \(2100\) TL kalıyor. Maaş?
\item Bir sayının \(\frac{3}{5}\)'inin \(2\) fazlası, \(\frac{1}{2}\)'sinin \(5\) fazlasına eşit. Sayı? \(\star\)
\end{enumerate}
\column{0.34\textwidth}
\pause
\textbf{Cevaplar:}
\begin{enumerate}
\item \(27\) \ (sayı \(63\))
\item \(120\) TL
\item \(225\) L
\item \(3600\) TL
\item \(30\)
\end{enumerate}
\end{columns}
\vspace{4pt}{\footnotesize \(\star\) işaretli soruların adım adım çözümü sonraki slaytlardadır.}
\end{frame}

% ---- Ders 2 alıştırma çözümleri ----
\soru{Ders 2 -- Alıştırma Çözümü (Soru 2)}
{Parasının \(\dfrac{1}{3}\)'ünü, sonra kalanın \(\dfrac{1}{4}\)'ünü harcayan birine \(60\) TL kalıyor. Başlangıçta kaç TL'si vardı?}
{Paydalar \(3\) ve \(4\) \(\Rightarrow\) para \(=12k\).\\[2pt]
\(\frac{1}{3}\)'ü: \(4k\) \(\to\) kalan \(8k\); \quad kalanın \(\frac{1}{4}\)'ü: \(2k\) \(\to\) kalan \(6k\).
\[6k=60\ \Rightarrow\ k=10\ \Rightarrow\ 12k=120\ \text{TL}\]
Sağlama: \(120\to 80\to 60\) \checkmark}

\soru{Ders 2 -- Alıştırma Çözümü (Soru 3)}
{Bir deponun \(\dfrac{2}{5}\)'i doludur. \(60\) litre su eklenince deponun \(\dfrac{2}{3}\)'ü doluyor. Deponun tamamı kaç litredir?}
{\[\frac{2}{3}-\frac{2}{5}=\frac{10-6}{15}=\frac{4}{15}\ \Rightarrow\ \frac{4}{15}V=60\ \Rightarrow\ V=60\cdot\frac{15}{4}=225\ \text{L}\]
Sağlama: \(\frac{2}{5}\cdot 225=90\), \ \(90+60=150=\frac{2}{3}\cdot 225\) \checkmark}

\soru{Ders 2 -- Alıştırma Çözümü (Soru 5)}
{Bir sayının \(\dfrac{3}{5}\)'inin \(2\) fazlası, aynı sayının \(\dfrac{1}{2}\)'sinin \(5\) fazlasına eşittir. Sayı kaçtır?}
{EKOK\((5,2)=10\): her terimi \(10\) ile çarpalım.
\[\frac{3x}{5}+2=\frac{x}{2}+5\ \Rightarrow\ 6x+20=5x+50\ \Rightarrow\ x=30\]
Sağlama: \(18+2=20=15+5\) \checkmark}


% =====================================================
\section{Ders 3: Yaş Problemleri}
\begin{frame}
\centering\Large\color{anamavi}\textbf{Ders 3}\\[6pt]\normalsize Yaş Problemleri: Yaş Farkı, Tablo Yöntemi, Yaş Oranı ve Ortalaması
\end{frame}

\subsection{Altın kurallar}
\begin{frame}{Yaş Problemleri: Temel Kurallar}
\begin{block}{Kural}
\begin{enumerate}
\item İki kişi arasındaki \textbf{yaş farkı hiç değişmez}.
\item \(t\) yıl sonra herkesin yaşı \(t\) artar; \(t\) yıl önce herkesin yaşı \(t\) eksiktir.
\item \(n\) kişinin yaşları toplamı \(t\) yıl sonra \(n\cdot t\) artar; yaş ortalaması \(t\) artar.
\item Yaşlar arasındaki \textbf{oran} (kat) zamanla değişir.
\end{enumerate}
\end{block}
\begin{alertblock}{Sık yapılan hatalar}
\begin{itemize}
\item ``\(6\) yıl sonra'' derken yalnızca bir kişiye \(6\) eklemek: herkese eklenir.
\item Bulunan yaşın anlamlı olduğunu kontrol edin: geçmişteki yaşlar negatif olamaz.
\end{itemize}
\end{alertblock}
\end{frame}

\soru{Yaş Toplamı -- Örnek 24}
{(a) Bir anne ile iki çocuğunun yaşları toplamı \(40\)'tır. \(5\) yıl sonra yaşları toplamı kaç olur?\\
(b) Ayşe ile Fatma'nın yaşları farkı \(10\)'dur. \(5\) yıl sonra farkları kaç olur?}
{(a) \(3\) kişi, her biri \(5\) yaş büyür: \ \(40+3\cdot 5=55\).\\[4pt]
(b) Yaş farkı değişmez: \(10\).}

\soru{Yaş -- Örnek 25}
{Bir baba \(40\), oğlu \(12\) yaşındadır. Kaç yıl sonra babanın yaşı oğlunun yaşının \(3\) katı olur?}
{\(t\) yıl sonra: baba \(40+t\), oğul \(12+t\).
\[40+t=3(12+t)\ \Rightarrow\ 40+t=36+3t\ \Rightarrow\ t=2\]
Sağlama: \(42=3\cdot 14\) \checkmark}

\begin{frame}{Tablo Yöntemi}
\begin{block}{Kural}
Satırlara kişileri, sütunlara zamanları yazın; her hücreyi tek bilinmeyenle doldurun.
\begin{center}
\renewcommand{\arraystretch}{1.25}
\begin{tabular}{lccc}
\toprule
 & \(t\) yıl önce & Şimdi & \(t\) yıl sonra\\\midrule
Kişi 1 & \(x-t\) & \(x\) & \(x+t\)\\
Kişi 2 & \(y-t\) & \(y\) & \(y+t\)\\\bottomrule
\end{tabular}
\end{center}
Denklemi, sorunun söylediği \textbf{sütundaki} ilişkiden kurun.
\end{block}
\end{frame}

\soru{Tablo Yöntemi -- Örnek 26}
{Bir annenin yaşı kızının yaşının \(4\) katıdır. \(6\) yıl sonra annenin yaşı kızının yaşının \(3\) katı olacaktır. Anne bugün kaç yaşındadır?}
{\begin{tabular}{lcc}
 & Şimdi & \(6\) yıl sonra\\\hline
Kız & \(x\) & \(x+6\)\\
Anne & \(4x\) & \(4x+6\)
\end{tabular}
\[4x+6=3(x+6)\ \Rightarrow\ x=12\ \Rightarrow\ \text{anne } 48\]
Sağlama: \(54=3\cdot 18\) \checkmark}

\soru{Tablo Yöntemi -- Örnek 27}
{Bir babanın yaşı oğlunun yaşının \(3\) katıdır. \(5\) yıl önce babanın yaşı oğlunun yaşının \(4\) katıydı. Baba bugün kaç yaşındadır?}
{\begin{tabular}{lcc}
 & \(5\) yıl önce & Şimdi\\\hline
Oğul & \(x-5\) & \(x\)\\
Baba & \(3x-5\) & \(3x\)
\end{tabular}
\[3x-5=4(x-5)\ \Rightarrow\ x=15\ \Rightarrow\ \text{baba } 45\]
Sağlama: \(40=4\cdot 10\) \checkmark}

\soru{Yaş Toplamı -- Örnek 28}
{Bir anne ile kızının yaşları toplamı \(48\)'dir. \(4\) yıl sonra annenin yaşı kızının yaşının \(3\) katı olacaktır. Anne bugün kaç yaşındadır?}
{Kız \(k\), anne \(48-k\). \ \(4\) yıl sonra: \(52-k\) ve \(k+4\).
\[52-k=3(k+4)\ \Rightarrow\ 40=4k\ \Rightarrow\ k=10,\quad \text{anne}=38\]
Sağlama: \(42=3\cdot 14\) \checkmark}

\subsection{Ben senin yaşındayken}
\begin{frame}{``Ben Senin Yaşındayken'' Tipi}
\begin{block}{Kural}
\(A\) büyük, \(B\) küçük olsun; yaş farkı \(d=A-B\).
\begin{itemize}
\item ``\(A\), \(B\)'nin yaşındayken'' \(=\) \(d\) yıl önce \(\Rightarrow\) o zaman \(B\)'nin yaşı \(B-d\).
\item ``\(B\), \(A\)'nın yaşına geldiğinde'' \(=\) \(d\) yıl sonra \(\Rightarrow\) o zaman \(A\)'nın yaşı \(A+d\).
\end{itemize}
Sayı doğrusunda: \(B-d,\ \ B,\ \ A,\ \ A+d\) eşit aralıklıdır (aralık \(d\)).
\end{block}
\begin{alertblock}{Sık yapılan hatalar}
``Ben senin yaşındayken'' ifadesindeki geçen süre, iki kişinin yaş farkıdır; ayrıca bir \(t\) aramayın.
\end{alertblock}
\end{frame}

\soru{Ben Senin Yaşındayken -- Örnek 29}
{Ali bugün \(24\) yaşındadır. Ali, Veli'nin bugünkü yaşındayken Veli \(12\) yaşındaydı. Veli bugün kaç yaşındadır?}
{Veli \(x\); fark \(d=24-x\). \ \(d\) yıl önce Veli \(x-d=12\):
\[24-x=x-12\ \Rightarrow\ 2x=36\ \Rightarrow\ x=18\]
Sağlama: \(6\) yıl önce Ali \(18\), Veli \(12\) \checkmark \quad (\(12,\ 18,\ 24\): aralık \(6\))}

\soru{Ben Senin Yaşındayken -- Örnek 30}
{Ali, Veli'nin bugünkü yaşındayken Veli \(10\) yaşındaydı. Veli, Ali'nin bugünkü yaşına geldiğinde Ali \(34\) yaşında olacaktır. İkisinin bugünkü yaşları kaçtır?}
{Sayı doğrusu: \(10,\ V,\ A,\ 34\) eşit aralıklı \(\Rightarrow\) aralık \(d\):
\[10+3d=34\ \Rightarrow\ d=8\ \Rightarrow\ V=18,\ A=26\]
Sağlama: \(8\) yıl önce \(18\) ve \(10\); \(8\) yıl sonra \(26\) ve \(34\) \checkmark}

\subsection{Oran, ortalama, doğum yılı}
\soru{Yaş Oranı -- Örnek 31}
{İki kardeşin yaşları oranı \(\dfrac{2}{3}\)'tür. \(4\) yıl sonra bu oran \(\dfrac{3}{4}\) olacaktır. Kardeşler bugün kaç yaşındadır?}
{Yaşlar \(2k\) ve \(3k\):
\[\frac{2k+4}{3k+4}=\frac{3}{4}\ \Rightarrow\ 8k+16=9k+12\ \Rightarrow\ k=4\]
Kardeşler \(8\) ve \(12\) yaşında. \quad Sağlama: \(\dfrac{12}{16}=\dfrac{3}{4}\) \checkmark}

\soru{Yaş Ortalaması -- Örnek 32}
{\(4\) kişilik bir ailenin bugünkü yaş ortalaması \(22\)'dir. \(3\) yıl sonra bu ailenin yaşları toplamı kaç olur?}
{Yol 1: toplam \(22\cdot 4=88\); \ artış \(4\cdot 3=12\) \(\Rightarrow\) \(100\).\\[4pt]
Yol 2: ortalama \(3\) artar: \(25\) \(\Rightarrow\) toplam \(25\cdot 4=100\).}

\soru{Yaş Farkı -- Örnek 33}
{Aslı, Kerem'den \(4\) yaş büyüktür. Kerem, Aslı'nın bugünkü yaşına geldiğinde ikisinin yaşları toplamı \(48\) olacaktır. Kerem bugün kaç yaşındadır?}
{Kerem \(k\), Aslı \(k+4\). Kerem'in Aslı'nın yaşına gelmesi \(4\) yıl sürer:
\[(k+4)+(k+8)=48\ \Rightarrow\ 2k=36\ \Rightarrow\ k=18\]
Sağlama: \(4\) yıl sonra \(22+26=48\) \checkmark}

\soru{Doğum Yılı -- Örnek 34}
{Elif 1998'de doğmuştur; annesi Elif doğduğunda \(27\) yaşındaydı. (a) Elif 2026'da kaç yaşına girer? (b) Annesi hangi yıl \(60\) yaşına girer? (c) Hangi yıl annenin yaşı Elif'in yaşının \(2\) katı olur?}
{(a) \(2026-1998=28\). \quad (b) Anne 1971 doğumlu: \(1971+60=2031\).\\[3pt]
(c) Fark \(27\) sabit: \(2e=e+27\Rightarrow e=27\) \(\Rightarrow\) \(1998+27=2025\).\\[3pt]
{\normalsize (Yaş \(=\) yıl farkı: doğum günü geçmiş kabul edilir.)}}

\begin{frame}{Ders 3 -- Sınıfta Siz Çözün}
\begin{columns}[t]
\column{0.62\textwidth}
\begin{enumerate}
\item Baba \(42\), kızı \(12\) yaşında. Kaç yıl sonra babanın yaşı kızının \(2\) katı olur? \(\star\)
\item \(5\) kişilik ailenin yaşları toplamı \(96\). \(4\) yıl sonra toplam?
\item Can \(30\) yaşında. Can, Ece'nin bugünkü yaşındayken Ece \(14\) yaşındaydı. Ece kaç yaşında? \(\star\)
\item İki kardeşin yaş oranı \(\frac{3}{4}\); \(5\) yıl önce \(\frac{2}{3}\) idi. Büyük kardeş kaç yaşında? \(\star\)
\item 2009 doğumlu biri hangi yıl \(40\) yaşına girer?
\end{enumerate}
\column{0.34\textwidth}
\pause
\textbf{Cevaplar:}
\begin{enumerate}
\item \(18\) yıl
\item \(116\)
\item \(22\)
\item \(20\) \ (\(15\) ve \(20\))
\item 2049
\end{enumerate}
\end{columns}
\vspace{4pt}{\footnotesize \(\star\) işaretli soruların adım adım çözümü sonraki slaytlardadır.}
\end{frame}

% ---- Ders 3 alıştırma çözümleri ----
\soru{Ders 3 -- Alıştırma Çözümü (Soru 1)}
{Bir baba \(42\), kızı \(12\) yaşındadır. Kaç yıl sonra babanın yaşı kızının yaşının \(2\) katı olur?}
{\(t\) yıl sonra: baba \(42+t\), kız \(12+t\).
\[42+t=2(12+t)\ \Rightarrow\ 42+t=24+2t\ \Rightarrow\ t=18\]
Sağlama: baba \(60\), kız \(30\) \checkmark \quad (Kısa yol: fark \(30\) sabit; kız \(30\) olmalı.)}

\soru{Ders 3 -- Alıştırma Çözümü (Soru 3)}
{Can \(30\) yaşındadır. Can, Ece'nin bugünkü yaşındayken Ece \(14\) yaşındaydı. Ece bugün kaç yaşındadır?}
{Ece \(x\), fark \(d=30-x\). \ \(d\) yıl önce Can \(x\) yaşındaydı, Ece \(x-d=14\).
\[30-x=x-14\ \Rightarrow\ 2x=44\ \Rightarrow\ x=22\]
Sağlama: \(14,\ 22,\ 30\) eşit aralıklı (\(d=8\)) \checkmark}

\soru{Ders 3 -- Alıştırma Çözümü (Soru 4)}
{İki kardeşin yaşları oranı \(\dfrac{3}{4}\)'tür. \(5\) yıl önce bu oran \(\dfrac{2}{3}\) idi. Büyük kardeş bugün kaç yaşındadır?}
{Bugün \(3k\) ve \(4k\); \(5\) yıl önce \(3k-5\) ve \(4k-5\).
\[\frac{3k-5}{4k-5}=\frac{2}{3}\ \Rightarrow\ 9k-15=8k-10\ \Rightarrow\ k=5\]
Kardeşler \(15\) ve \(20\) \(\Rightarrow\) büyük \(20\). \ Sağlama: \(\dfrac{10}{15}=\dfrac{2}{3}\) \checkmark}


\begin{frame}{Özet}
\begin{columns}[t]
\column{0.32\textwidth}
\begin{block}{Ders 1}
\begin{itemize}
\item ``\(5\) eksiğinin \(3\) katı'' \(=3(x-5)\)
\item Ardışık tek/çift: fark \(2\)
\item \(\ab{ab}-\ab{ba}=9(a-b)\), \ \(\ab{ab}+\ab{ba}=11(a+b)\)
\end{itemize}
\end{block}
\column{0.32\textwidth}
\begin{block}{Ders 2}
\begin{itemize}
\item Bütün \(=\) paydaların EKOK'u kadar birim
\item Kalanın kesri: önce kalanı bul
\item Depo: kesir farkı \(=\) eklenen litre
\end{itemize}
\end{block}
\column{0.32\textwidth}
\begin{block}{Ders 3}
\begin{itemize}
\item Yaş farkı sabit
\item \(n\) kişi, \(t\) yıl: toplam \(+\,n\cdot t\)
\item Tablo: Önce / Şimdi / Sonra
\end{itemize}
\end{block}
\end{columns}
\vspace{6pt}
\begin{alertblock}{Gelecek hafta}
Yüzde, kâr--zarar, karışım ve faiz problemleri; genel tekrar
\end{alertblock}
\end{frame}

\begin{frame}
\centering\Large İyi çalışmalar
\end{frame}

\end{document}
